% !TEX root = ../../tree_notes.tex
% =====================================================================
%  7.1  The Binary Heap                          sec:c07-heap
%  Exports: def:c07-heap  lem:c07-height  def:c07-sift
%           lst:c07-heapops  ex:c07-sift-run  (fig:c07-layout, fig:c07-heapR
%           are defined in 00_opener.tex and reused here)
%           prob:c07-stone  lst:c07-stone-soln   (practice: LC 703)
% =====================================================================
\section{The Binary Heap}\label{sec:c07-heap}

\begin{previouslybox}
\begin{enumerate}[nosep, leftmargin=1.6em]
  \item A complete binary tree with $n$ nodes fills every level except possibly the last, and the last
        from the left. Why can such a tree be stored with \emph{no pointers} at all? (its shape is fixed,
        so position determines parent/child -- Chapter~4, \chref{c04})
  \item On a BST the \emph{inorder} walk lists keys sorted. What single promise does a heap make instead,
        and what does it deliberately \emph{not} promise between siblings? (parent beats children; siblings
        unordered -- \chref{c05})
  \item A \emph{min}-heap puts its smallest element where, relative to a \emph{max}-heap's largest? (both
        at the root; the rule is just flipped -- the read-only cousin is Chapter~9, \chref{c26})
  \item How many elements can sit on the bottom level of a tree of height $h$, and how tall is a tree of
        $n$ nodes if it is as short as shape allows? (up to $2^{h}$; height $\approx\lg n$)
\end{enumerate}
\end{previouslybox}

\begin{topicintro}
Chapter~9 (\chref{c26}) \emph{read} a heap order straight out of a frozen array. This section starts the
moving version: a heap you \emph{edit}. The whole data structure rests on one deliberately weak promise
-- \textbf{the root is the extreme, and that is all}. A \term{binary heap} is a complete binary tree kept
in a plain array and held to the \term{heap order} ``every parent $\le$ its children'' (min-heap) or the
reverse (max-heap). Crucially there is \emph{no} left-to-right order between siblings: the array position
$t_3$ may hold a smaller value than $t_2$ with no rule broken. That weak promise is exactly what a greedy
or a top-$k$ question ever asks for -- ``give me the current extreme now'' -- and it is shockingly cheap
to maintain, because the tree's shape is fixed by the array index and never needs a rotation.

Two moves do everything. \term{Sift-up} plants a new leaf at the end and bubbles it toward the root while
it beats its parent; \term{sift-down} drops the last element onto the vacated root and pushes it down
toward the leaves while it is beaten by a child. Both walk a single root--leaf path, and the heap is
\emph{short by shape}: its height is $\lfloor\lg n\rfloor$ (\lemref{c07-height}). So \code{push} and
\code{pop} are each $\bigO{\log n}$, and \code{top} -- read the extreme -- is one array cell. We fix the
layout, define the two sifts, hand-write the three operations, and spend them on LeetCode 1046, the
purest ``extract-the-extreme-then-reinsert'' warm-up in the book.
\textbf{Prerequisites.} Array index arithmetic; the complete-binary-tree shape previewed in Chapter~4
(\chref{c04}). No search-tree ordering is assumed; indeed we open by contrasting the heap's ``no sibling
order'' against the BST inorder of Chapter~7 (\chref{c05}).

\textbf{After this section you can}
\begin{itemize}[nosep]
  \item store a heap in an array, compute parent/children from an index, and state exactly what the heap
        order promises and what it leaves unordered; \hfill\lv{2}--\lv{3}
  \item write sift-up and sift-down push/pop cold, and explain why each is $\bigO{\log n}$ from the height
        bound; \hfill\lv{3}
  \item recognise a ``repeatedly grab the extreme'' problem and reach for \code{priority\_queue} rather
        than a sort; \hfill\lv{3}
\end{itemize}
\end{topicintro}

\begin{recognitionbox}[Pattern card -- the extreme of a \emph{changing} set]
\reckey{you must repeatedly read (and remove) the \emph{current} smallest or largest element of a
multiset that you also insert into, and you never ask for the second-smallest, a predecessor, or an
arbitrary rank -- only the one element at the top.}
\recuse{
  \item ``always smash the two heaviest'' / ``combine the two smallest'' -- pop the extreme, push a
        result back, \probref{c07-stone} and \probref{c07-connect};
  \item ``the $k$-th largest so far'' / ``keep the top $k$'' -- a size-$k$ heap, \secref{c07-build};
  \item ``schedule by priority'', ``which task runs next'' -- pop the highest-priority item, the graph
        chapters' Dijkstra/Prim engine (previewed, not built here).
}
\recnotuse{
  \item you need \emph{any} rank, predecessor, or successor -- that is a total order, go to the BST/treap
        of Chapters~7--8 (\chref{c05}, \chref{c06}); a heap cannot answer ``what is just below $x$?'';
  \item the set is \emph{static} and you only want its minimum -- one scan, no structure, or the read-only
        Cartesian tree of Chapter~9 (\chref{c26});
  \item you need the extreme of \emph{every subarray} $[l,r]$ -- that is range-min, a sparse table or
        segment tree (\chref{c13}, \chref{c26}).
}
One-line test: \emph{does the question keep asking ``who is the biggest/smallest \emph{right now}'' after
each edit? Then a heap is the tool, and the two sifts below are all the machinery you need.}
\end{recognitionbox}

\subsection{Intuition}

Lay the running array $R=(5,3,8,1,4,2,7)$ into the slots of a complete tree, numbering the slots from $1$
(\figref{c07-layout}). Nothing about the values has been changed -- it is the same seven numbers -- but now
each slot has a fixed address, and the address encodes the tree: slot $i$'s parent is $\lfloor i/2\rfloor$,
its children are $2i$ and $2i+1$. No pointer is stored anywhere; the tree is \emph{implicit in the index}.
That is the entire reason a heap is a handful of integers and one loop: the shape (complete, filled from
the left) is forced, so the only thing you ever recompute is \emph{where a value sits}, never the
structure itself.

Now impose the heap order. Read $R$ as a min-heap and you must have each parent $\le$ its children; that
is \emph{not} true of $R$ as written ($5$ sits above $3$ and $8$). But once we heapify it -- by the sifts
below, and $\bigO{n}$ at build time in \secref{c07-build} -- we get the array $(1,3,2,5,4,8,7)$ shown as
\figref{c07-heapR}. Look hard at its root: $1$, the minimum, sits at index~$1$. Look at the two children
of the root: $3$ on the left, $2$ on the right. The \emph{left} child is \emph{bigger} than the right
child. A BST would consider that a corruption; a heap does not care, because it makes \emph{no} promise
across siblings. And in fact the whole right half $\{2,8,7\}$ is smaller than much of the left half, yet
the structure is perfectly valid. ``The top is the extreme; everything below is only loosely organised''
is the deal, and it is precisely loose enough to be an array and tight enough to answer the one question
greedy problems ask.

\begin{figure}[H]
\centering
\begin{tikzpicture}[diagram, x=13mm, y=9mm, level distance=13mm, sibling distance=22mm]
  \node[troot] (r) {$1$}
    child { node[tnode] {$3$}
      child { node[tleaf] {$5$} }
      child { node[tleaf] {$4$} } }
    child { node[tnode] (c2) {$2$}
      child { node[tleaf] {$8$} }
      child { node[tleaf] {$7$} } };
  \draw[tedgehl] (r) -- (c2);
  \node[annot, align=left, font=\scriptsize] at (-4.3,0.5)
    {\textbf{root $=1$} is the\\ minimum of $R$\\ (\emph{the} promise)};
  \node[annot, align=left, font=\scriptsize] at (4.2,-1.2)
    {left child $3 >$ right child $2$\\ -- \emph{allowed}: no\\ sibling order in a heap};
\end{tikzpicture}
\caption*{The min-heap $(1,3,2,5,4,8,7)$, i.e.\ \figref{c07-heapR}, annotated. The highlighted root--right
edge exists only to draw your eye to the fact that the value grows \emph{down} both subtrees ($1\le3\le5$,
$1\le2\le8$) while the left and right branches are mutually unordered. That asymmetry is why a heap cannot
find the \emph{second} smallest without a second pop, and why it costs so little to keep.}
\end{figure}

\subsection{Definitions and Notation}

\begin{definition}{Binary heap (array representation)}{c07-heap}
A \term{(min-)binary heap} is a complete binary tree stored in an array $h[1..n]$ satisfying the
\term{heap-order} property: for every node $i>1$, $h[\lfloor i/2\rfloor] \le h[i]$. Equivalently, every
parent is $\le$ each of its children. A \term{max-heap} reverses the inequality. No other ordering is
implied: for two nodes that are not in an ancestor--descendant relation there is \emph{no} guarantee, so
in particular \text{left}(i) and \text{right}(i) may hold values in either order.
\end{definition}

\noindent\textbf{In plain words.} It is a complete tree, so it lives in an array with parent
$\lfloor i/2\rfloor$ and children $2i,2i+1$. It is a \emph{heap}, so each node is $\le$ (or $\ge$) its
children -- the extreme floats to index~$1$, and \emph{nothing else} is sorted.

\begin{lemma}{Height of a heap is $\lfloor\lg n\rfloor$}{c07-height}
A heap on $n\ge 1$ elements has height $\lfloor\lg n\rfloor$ (counting edges from root to the deepest
leaf). Hence any root-to-leaf path -- and every sift -- has length $O(\log n)$.
\end{lemma}
\begin{proofbox}
A complete tree of height $h$ has its levels $0,1,\dots,h-1$ full (that is $1+2+\cdots+2^{h-1}=2^{h}-1$
nodes) and its last level non-empty, so $2^{h}-1 < n \le 2^{h+1}-1$, i.e.\ $2^{h}\le n < 2^{h+1}$. Taking
$\lg$ gives $h \le \lg n < h+1$, so $h=\lfloor\lg n\rfloor$. The filling-from-the-left rule is what forces
every level but the last to be full; because a heap \emph{never rotates} -- it only adds or removes the
last slot -- the shape is guaranteed after every edit, and so is the bound. Every sift walks at most one
such path, hence $\bigO{\log n}$.
\end{proofbox}

\begin{definition}{Sift-up and sift-down}{c07-sift}
\term{Sift-up}($i$) repeatedly compares $h[i]$ with its parent $\lfloor i/2\rfloor$ and swaps them while
$i>1$ and $h[i] < h[\lfloor i/2\rfloor]$ (min-heap), restoring the heap order up the single path from a
freshly added leaf. \term{Sift-down}($i$) repeatedly swaps $h[i]$ with its \emph{smallest} child while
that child is $<h[i]$ (min-heap), restoring the heap order down the single path from the root. A sift is
correct precisely because only one path's heap-order edge can be violated after a \code{push} or \code{pop}
-- every untouched subtree stays a heap -- so fixing that one path repairs the whole structure.
\end{definition}

\begin{insightbox}
The reason a heap is $\bigO{\log n}$ and a \emph{sorted array} is $\bigO{n}$ per insertion is the same
reason it is not a BST: the heap refuses to sort. A sorted array must move every element to open a gap at
the right rank; a heap only needs the extreme to be reachable at index~$1$, so it can let a value land
anywhere in $\bigO{\log n}$ swaps along a root-to-leaf path and call it done. You are paying the minimum
possible price for the \emph{one} guarantee the problem actually needs. When you later need richer order,
you graduate to the ordered trees of Chapters~7--8 (\chref{c05}, \chref{c06}).
\end{insightbox}

\subsection{The Algorithm: Two Sifts, Three Operations}

Everything below is a $\bigO{\log n}$ walk down or up \emph{one} path of the implicit array tree. We keep
$h$ $1$-based with $h[0]$ an unused dummy so the parent/child formulas are exactly $\lfloor i/2\rfloor$ and
$2i,2i+1$ -- matching \defref{c07-heap} and every trace in this chapter.

\begin{lstlisting}[caption={A hand-written min-heap: \code{push}, \code{pop}, \code{top} on an implicit array}, label={lst:c07-heapops}]
vector<int> h(1, 0);                 // 1-based min-heap; h[0] is an unused dummy
int parent(int i){ return i >> 1; }          // floor(i/2)
int left    (int i){ return i << 1; }        // 2i
int right   (int i){ return (i << 1) | 1; }  // 2i+1

void push(int x) {                   // plant a new leaf, then sift UP
    h.push_back(x);
    int i = (int)h.size() - 1;
    while (i > 1 && h[i] < h[parent(i)]) {
        swap(h[i], h[parent(i)]);
        i = parent(i);
    }
}
int pop() {                          // remove min: refill root, sift DOWN
    int res = h[1];
    h[1] = h.back();                 // last element takes the root's place
    h.pop_back();
    int i = 1, n = (int)h.size() - 1;
    while (true) {
        int s = i;                   // candidate to swap with = smallest child
        if (left(i)  <= n && h[left(i)]  < h[s]) s = left(i);
        if (right(i) <= n && h[right(i)] < h[s]) s = right(i);
        if (s == i) break;           // already <= both children: heap restored
        swap(h[i], h[s]);
        i = s;
    }
    return res;
}
int top() { return h[1]; }           // the extreme is a single array cell away
\end{lstlisting}

\begin{complexitybox}
Time: \code{push} sifts up at most one root--leaf path, \code{pop} sifts down one path -- by
\lemref{c07-height} each is $O(\height)=\bigO{\log n}$ swaps; \code{top} and a ``is it empty'' test are
$\bigO{1}$; building by pushing $n$ values one at a time is $\bigO{n\log n}$ (\secref{c07-build} beats
that to $\bigO{n}$). Memory $\bigO{n}$ for the array -- and \emph{nothing else}: no pointers, no recursion
stack. Verdict: this is the smallest non-trivial priority queue you can write; keep it in memory so that,
when a contest forbids the library, you type it in seconds.
\end{complexitybox}

\subsection{Worked Example}

\begin{example}{Sift-up build and one sift-down pop, on $R=(5,3,8,1,4,2,7)$}{c07-sift-run}
\emph{Sift-up (build by inserting in array order).}\\
$5$: $[5]$. \ $3$: append, parent $5>3$ swap $\to[3,5]$. \ $8$: parent $3\le8$, stop $\to[3,5,8]$.\\
$1$: append at index $4$; parent $h[2]{=}5>1$ swap, then parent $h[1]{=}3>1$ swap $\to[1,3,8,5]$.\\
$4$: index $5$, parent $h[2]{=}3\le4$, stop $\to[1,3,8,5,4]$.\\
$2$: index $6$, parent $h[3]{=}8>2$ swap $\to[1,3,2,5,4,8]$; now parent $h[1]{=}1\le2$, stop.\\
$7$: index $7$, parent $h[3]{=}2\le7$, stop $\to[1,3,2,5,4,8,7]$.\\
Result $(1,3,2,5,4,8,7)$ -- exactly \figref{c07-heapR}. Note index $2{=}3$ and index $3{=}2$ are
\emph{out of order} and the heap is still valid.\\[2pt]
\emph{Sift-down (one \code{pop}, removing the min $1$).}\\
Move last element $7$ to root, shrink: $[7,3,2,5,4,8]$. Sift $7$ down: children $h[2]{=}3$, $h[3]{=}2$,
smallest child $2$ at index $3$; $7>2$ swap $\to[2,3,7,5,4,8]$. Now at index $3$, its only in-range child
is $h[6]{=}8$; $7\le8$, stop. Heap $[2,3,7,5,4,8]$, and the popped value was $1$. \checkmark\ The new root
$2$ is the second-smallest of $R$: the heap hands you the extremes \emph{in order}, one \code{pop} at a
time.
\end{example}
\begin{tracebox}[Array view of \exref{c07-sift-run}, indices 1-based]
after build: i=1..7 hold 1 3 2 5 4 8 7 (root 1; children of 1 are 3,2; of 3 are 5,4; of 2 are 8,7).\\
pop(): root$\leftarrow$7 (last), size 6. path: 7@1 vs min-child 2@3 -> swap; 7@3 vs child 8@6 -> stop.\\
final: 2 3 7 5 4 8. Every parent $\le$ children: 2<=3,2; 3<=5,4; 7<=8. Sibling 3>2 -- allowed.
\end{tracebox}


\subsection{Variations and Twists}

\begin{itemize}[nosep]
  \item \textbf{Min vs.\ max}: flip every comparison in \coderef{c07-heapops}; a max-heap is the same code
        with $>$ where the min-heap has $<$. Which one you want is decided by which extreme the problem
        greedily removes. \hfill\lv{2}
  \item \textbf{Arbitrary delete-min of a non-root}: only the \emph{root} is guaranteed extreme. To remove
        a specific interior element you must either sift it out after swapping in the last leaf (breaking
        the ``complete'' shape -- then repair) or mark it dead and lazily pop stale tops; the lazy trick
        powers the Summit (\probref{c07-summit}). \hfill\lv{4}
  \item \textbf{Decrease-key / increase-key}: a plain binary heap has \emph{no} fast handle for ``element
        at unknown index changed'' -- you need a position map (index $\leftrightarrow$ handle), the thing
        Dijkstra's heap variant carries. Not free here; noted for the graph chapters. \hfill\lv{4}
  \item \textbf{k-th smallest}: popping $k$ times from a min-heap is $\bigO{k\log n}$; for $k\ll n$ a
        size-$k$ \emph{max}-heap is better (\secref{c07-build}). The heap gives extremes \emph{in order},
        not ranked access. \hfill\lv{3}
\end{itemize}

\subsection{Common Pitfalls}

\begin{warningbox}[Trap 1 -- assuming the array is sorted]
Reading a heap's array left-to-right gives $1,3,2,5,4,8,7$: it is \emph{not} sorted, and no slice of it is.
The only invariant is ``parent beats children'' (\defref{c07-heap}); siblings are free. Trying to binary
search a heap, or expecting $h[i]\le h[i+1]$, is the classic wrong-structure error. If you need sorted
order, either pop everything ($\bigO{n\log n}$, the basis of heapsort) or use the ordered trees of
Chapters~7--8 (\chref{c05}, \chref{c06}).
\end{warningbox}
\begin{lstlisting}[caption={\textcolor{warningred}{WRONG}: treating a heap array as sorted}, language=C++, basicstyle=\ttfamily\small]
// h is a min-heap array; this is NOT a valid way to find the 2nd smallest:
int secondSmallest = h[1];   // h[1] is the min; h[2] is a child of root, not "2nd smallest"
\end{lstlisting}
\begin{lstlisting}[caption={\textcolor{successgreen}{RIGHT}: the 2nd extreme comes from a second pop}, language=C++, basicstyle=\ttfamily\small]
int mn = pop();              // remove the min; heap re-sifts, O(log n)
int second = top();          // now the root is the 2nd smallest
push(mn);                    // (restore if you must keep the element)
\end{lstlisting}

\begin{warningbox}[Trap 2 -- off-by-one in the index formulas]
The layout only works when the tree is \emph{complete and filled from the left}. Use $1$-based indexing
(parent $i/2$, children $2i,2i{+}1$) or $0$-based (parent $(i-1)/2$, children $2i{+}1,2i{+}2$) -- but do
\emph{not} mix them. Every bug in \coderef{c07-heapops} that ``sometimes leaves a parent bigger than a
child'' is a mixed-base index. We standardise on $1$-based throughout this chapter.
\end{warningbox}

\begin{warningbox}[Trap 3 -- sifting the wrong way after \code{pop}]
After \code{pop} you replace the root with the \emph{last leaf} and sift \emph{down}; after \code{push} you
append at the end and sift \emph{up}. Sifting the new root \emph{up} (it has no parent) is a no-op that
silently leaves the heap broken at the bottom. The direction is forced by which single edge you broke:
\code{push} breaks the edge to a new leaf (fix upward); \code{pop} breaks the root's edges (fix downward).
\end{warningbox}

\subsection{Problem Lab}

The section's anchor is LeetCode 1046, \probref{c07-stone}: a pile of stones, repeatedly smash the two
heaviest. It is the most literal possible ``extract the extreme, then maybe reinsert a residue'' loop --
exactly the \code{pop}$\to$\code{pop}$\to$\code{push} rhythm \coderef{c07-heapops} provides -- so it is the
first problem every reader should be able to write cold with a \code{priority\_queue}.

% ---------------------------------------------------------------------------
\solvedproblem{c07-stone}{Last Stone Weight}{LeetCode 1046}
  {https://leetcode.com/problems/last-stone-weight/}{LC Easy · L2}
  {max-heap greedy extraction, pop-pops-push loop, priority\_queue}

\begin{problemstatement}[Last Stone Weight]
You are given an array \texttt{stones} where \texttt{stones[i]} is the weight of the $i$-th stone. On each
turn you choose the two \textbf{heaviest} stones with weights $x \le y$ and smash them: if $x = y$ both are
destroyed; otherwise the stone of weight $x$ is destroyed and the stone of weight $y$ becomes weight
$y - x$. Return the \textbf{smallest possible weight} of the last remaining stone, or $0$ if none remain.
Constraints: $1 \le \texttt{stones.length} \le 30$ and $1 \le \texttt{stones[i]} \le 1000$.
\end{problemstatement}

\begin{attemptgate}[Attempt it yourself first -- \lv{2}, give yourself 10 minutes]
``Always take the two heaviest'' is not a sorting problem; it is a ``pop the extreme twice'' problem. Model
one turn directly as two \code{pop}s and at most one \code{push}.
\peek{What container lets you get the current maximum in $\bigO{\log n}$ and then remove it?}
\nudge{Use a \emph{max}-heap. Each turn: pop the top twice (the two heaviest $y$ and $x$), and if
$y \ne x$ push $y-x$ back. Repeat while the heap has at least two stones.}
\fullhint{When the loop ends either the heap is empty (answer $0$) or one stone remains (answer that
stone, the top). For the sample $[2,7,4,1,8,1]$ the sequence of heaps is: smash $8,7\to1$; $4,2\to2$;
$2,1\to1$; $1,1\to$ both gone; leaving a single $1$. Answer $=1$.}
\end{attemptgate}

\solstep{Understanding the Problem}
We repeatedly act on the \emph{two largest} elements, replacing them by their (signed, non-negative)
difference, until fewer than two stones remain. The result depends only on \emph{which} two are largest at
each step -- a pure greedy on the current extremes. Because each smash can create a new value ($y-x$) that
is again among the largest, the extreme must be recomputed every turn; sorting once is wrong. That is the
signature of a priority queue: extract, mutate, reinsert, extract again.

\solstep{Key Observations}
\begin{observation}
A turn touches only the two heaviest stones, so a \emph{max}-heap gives them in $\bigO{\log n}$ each;
everything below the top is irrelevant to the decision, which is exactly the weak promise a heap makes
(\defref{c07-heap}).
\end{observation}
\begin{observation}
After popping $y\ge x$: push $y-x$ iff $y>x$ (if equal, both are destroyed and nothing is pushed). The
multiset size drops by $1$ or $2$ per turn, so at most $n-1$ turns occur.
\end{observation}
\begin{observation}
The process is deterministic given the multiset (the two heaviest are well-defined), so the greedy ``answer
is whatever the loop leaves'' is automatically optimal -- there is no choice to make.
\end{observation}

\solstep{Think It Through With Me}
\begin{thinkbox}
\deadend{Sort the array once and always look at the last two entries.}
{The smash produces $y-x$ which may belong anywhere in the order; you would re-sort every turn --
$\bigO{n^2\log n}$ and easy to get wrong. Re-sorting is the tell you wanted a heap.}
\deadend{Scan for the two maxima each turn.}
{$\bigO{n}$ per turn, $\bigO{n^2}$ total; fine for $n\le 30$ but misses the \emph{point}: the operation is
``extreme of a changing set,'' which is a priority queue, not a scan.}
\aha{Pop the top twice, push their difference back when non-zero. A max-heap keeps the two heaviest on tap
in $\bigO{\log n}$, and the loop is literally the \code{pop}$\to$\code{pop}$\to$\code{push} rhythm of
\coderef{c07-heapops}.}
\end{thinkbox}

\solstep{Solution}
Load all stones into a max-heap. While its size is at least $2$: pop $y$ (largest), pop $x$ (next largest);
if $y \ne x$ push $y-x$. When the loop ends, return $0$ if the heap is empty, else return the single
remaining top. Correctness is immediate from the greedy; the only implementation care is that the standard
\code{priority\_queue} is a \emph{max}-heap by default, which is exactly what we want here.

\solstep{Proof of Correctness}
\begin{proofbox}[The loop returns the smallest possible last weight]
Each turn is fully determined by the current multiset: the two heaviest stones $x\le y$ are unique (as
\emph{values}), and the mandated outcome -- remove both, reinsert $y-x$ if $y>x$ -- matches the problem's
smash rule exactly, so the loop reproduces the prescribed process step for step. The loop invariant ``the
heap holds precisely the multiset of stones still standing'' holds initially and is preserved by each turn.
The loop stops iff fewer than two stones remain, i.e.\ the terminal states the problem enumerates: empty
(return $0$) or a single stone (return it). Since the problem's own process is deterministic, the weight
the loop returns \emph{is} the (uniquely determined, hence smallest-possible) answer. \end{proofbox}

\solstep{Implementation}
\begin{lstlisting}[caption={Last Stone Weight -- two pops and a conditional push on a max-heap (LeetCode 1046)}, label={lst:c07-stone-soln}]
class Solution {
public:
    int lastStoneWeight(vector<int>& stones) {
        priority_queue<int> pq(stones.begin(), stones.end()); // max-heap
        while (pq.size() > 1) {
            int y = pq.top(); pq.pop();   // heaviest
            int x = pq.top(); pq.pop();   // second heaviest
            if (x != y) pq.push(y - x);   // one shattered remnant
        }
        return pq.empty() ? 0 : pq.top(); // last standing stone
    }
};
\end{lstlisting}

\solstep{Trace}
\begin{tracebox}[Official sample: \texttt{stones = [2,7,4,1,8,1]} (the pile $S$)]
max-heap $\{8,7,4,2,1,1\}$. Turn 1: pop $8,7$, push $1$ $\to\{4,2,1,1,1\}$. Turn 2: pop $4,2$, push $2$
$\to\{2,1,1,1\}$. Turn 3: pop $2,1$, push $1$ $\to\{1,1,1\}$. Turn 4: pop $1,1$, equal so push nothing
$\to\{1\}$. Size $<2$: stop. One stone remains, top $=1$. \checkmark\ matches the official output $1$.
\end{tracebox}

\solstep{Complexity}
\begin{complexitybox}
Time: at most $n-1$ turns; each does two \code{pop}s and at most one \code{push}, each $\bigO{\log n}$, so
$\bigO{n\log n}$ overall (building the heap from the array is $\bigO{n}$, \secref{c07-build}). Memory
$\bigO{n}$. With $n\le 30$ anything passes, but the $\bigO{n\log n}$ heap is the \emph{transferable}
answer: it is the same loop that solves the Summit (\probref{c07-summit}) and every ``combine the two
extremes'' greedy. Verdict: \ac.
\end{complexitybox}

\solstep{Pitfalls and Variants}
\begin{warningbox}
Three ways to fall: (i) using a \emph{min}-heap -- the problem wants the two \emph{heaviest}, so you need a
max-heap (the default \code{priority\_queue}) or to negate every weight; (ii) forgetting the ``both
destroyed'' case and pushing $y-x=0$, which pollutes the heap with zero-weight phantom stones (harmless to
the answer but a size/logic trap when you port to $k$-way variants); (iii) writing the loop condition as
\code{size >= 1} instead of $>1$, then popping the top a second time from a one-element heap -- \re.
\emph{Variant:} \emph{Last Stone Weight II} (LC 1049) asks for the \emph{minimum} final weight over
\emph{any} pairing and is a subset-sum DP, \emph{not} a greedy -- a sharp reminder that ``smallest possible
over all orderings'' is a different question from ``the greedy's result.''
\end{warningbox}

\begin{insightbox}[Key takeaways]
\begin{enumerate}[nosep, leftmargin=1.6em]
  \item A heap is a \emph{complete tree in an array} promising only ``parent $\le$ children''
        (\defref{c07-heap}); the root is the extreme, siblings are unordered -- this is what makes it cheap
        and what makes it \emph{not} a sorted array or a BST.
  \item \code{push} sifts \emph{up} a new leaf, \code{pop} refills the root with the last leaf and sifts
        \emph{down}; both walk one path of length $\lfloor\lg n\rfloor$ (\lemref{c07-height},
        \defref{c07-sift}), so every operation is $\bigO{\log n}$.
  \item ``Repeatedly grab the extreme, mutate, reinsert'' -- \probref{c07-stone} -- is the shape of the
        whole chapter; the Summit (\probref{c07-summit}) is the same loop with a subtler greedy rule.
\end{enumerate}
\end{insightbox}

\begin{hardeningbox}
\begin{itemize}[nosep, leftmargin=1.4em]
  \item \textbf{Write \coderef{c07-heapops} from memory}, then run \exref{c07-sift-run}'s insert sequence
        $5,3,8,1,4,2,7$ through your own code and check you land on $(1,3,2,5,4,8,7)$.
  \item \textbf{Negation trick}: make a min-heap out of a max-heap by pushing $-w$ for every weight; verify
        \probref{c07-stone} still returns $1$.
  \item \textbf{Read ahead}: the $\bigO{n}$ \emph{build} you skipped here (you pushed one at a time) is
        \secref{c07-build}'s Floyd heapify -- and the size-$k$ heap for ``top $k$'' is right after it.
\end{itemize}
\end{hardeningbox}

\subsection{Practice Problems}

\begin{stuckbox}
\begin{itemize}[nosep, leftmargin=1.4em]
  \item Did you pick the heap flavour (max vs.\ min) to match \emph{which} extreme the greedy removes?
  \item Did you make the loop stop when fewer than the needed number of elements remain, not when it is
        merely empty?
  \item When a smash/merge produces a new value, did you \emph{reinsert} it so it can be re-chosen?
\end{itemize}
\end{stuckbox}

\begin{practiceset}
  \practice{Kth Largest Element in a Stream}{LeetCode 703}
    {https://leetcode.com/problems/kth-largest-element-in-a-stream/}{LC Easy · L2}
    {Design a class that, given an initial array and a value $k$, supports \texttt{add(val)} returning the
    current $k$-th largest element after each insertion. A \emph{min}-heap capped at size $k$ is the whole
    idea: push each new value, and the moment the heap exceeds $k$ elements pop its root. The surviving
    $k$ elements are exactly the $k$ largest seen so far, and the \emph{root} (their minimum) is the
    $k$-th largest -- report \code{pq.top()} in $\bigO{1}$. Each \texttt{add} is $\bigO{\log k}$. Do
    \emph{not} keep all elements or a max-heap of everything: you only ever need the frontier of the top
    $k$. Build the initial heap in $\bigO{n}$ first (\secref{c07-build}) if the seed array is large. This
    is the running ``top-$k$ frontier'' pattern used again in \probref{c07-topk}.}
\end{practiceset}

\begin{recallbox}[10-second refresher -- 7.1 The Binary Heap]
A binary heap is a \emph{complete} tree stored in an array (parent $\lfloor i/2\rfloor$, children
$2i,2i{+}1$) held to ``parent $\le$ children'' (\defref{c07-heap}); the root is the extreme and
\emph{siblings are unordered} -- that weak promise is why it is only an array plus one loop. Its height is
$\lfloor\lg n\rfloor$ (\lemref{c07-height}), so \code{push} (sift-up a new leaf) and \code{pop} (refill the
root, sift-down) are each $\bigO{\log n}$ (\defref{c07-sift}, \coderef{c07-heapops}); the array is
\emph{not} sorted. LeetCode 1046 (\probref{c07-stone}) is the pure extract-extreme loop, and LC 703 the
size-$k$ frontier.
\selftest{
  \item What does a heap promise between a node and its children, and what does it promise between
        siblings?
  \item Why is a heap's array \emph{not} sorted, and how do you correctly get the second-smallest?
  \item State \lemref{c07-height} and explain why it makes every sift $\bigO{\log n}$.
  \item After \code{pop}, why do you sift \emph{down} and not \emph{up}?
  \item For ``two heaviest stones'' do you want a min- or a max-heap, and why?
}
\end{recallbox}
